\section[$\et(\Bun_G)$]{$\et(\Bun_G)$（12 月 21 日）}

\srcpage{84}

\subsection{行政事务}

下次课将在 1 月 8 日星期五进行。课程将持续到 2 月 12 日星期五。

\subsection{我们进行到哪一步了？}

设 $G/E$ 是约化群。

我们定义了 Fargues--Fontaine 曲线上 $G$-丛的模叠 $\Bun_G$。它表示如下函子：
\[
S\in\Perf_{\ol{\Fq}}\mapsto\{X_S\text{ 上的 }G\text{-丛}\}.
\]
我们证明了：

\begin{theorem}
\begin{enumerate}[label=(\arabic*)]
\item $\Bun_G$ 是 Artin v-叠，上同调光滑，维数为 $0$。
\item 映射 $\lvert\Bun_G\rvert\to B(G)$ 是连续双射。
\item 对任意 $b\in B(G)$，得到局部闭分层 $\lvert\Bun^b_G\rvert\subset\Bun_G$。它具有形式 $\Bun^b_G=[*/G_b]$，其中 $G_b$ 适合短正合列
\[
1\to\underbrace{\text{「幂幺群菱形」}}_{\text{正 Banach--Colmez 空间的迭代扩张}}\to G_b\to G_b(E)\to 1.
\]
选取代表元 $b\in G(\breveE)$ 诱导分裂 $[*/G_b(E)]\to[*/G_b]=\Bun^b_G$，它作为态射上同调光滑，相对维数为 $\langle2\rho,\nu_b\rangle$。
此外，$[*/G_b(E)]$ 上同调光滑，且同构于 $\Bun^1_{G_b}$，后者是维数为 $0$ 的 Artin v-叠。所以每个 $\Bun^b_G$ 也是上同调光滑的 Artin v-叠，维数为 $-\langle2\rho,\nu_b\rangle$。
\end{enumerate}
\end{theorem}

\begin{corollary}
Kottwitz 映射诱导 $\kappa\colon\pi_0\Bun_G\xrightarrow{\ \sim\ }\pi_1(G)_\Gamma$。等价地，$\Bun_G$ 的每个连通分支都是某个 $\Bun^b_G$ 的闭包，其中 $b\in B(G)$ 唯一且是基本的。
\end{corollary}

\begin{proof}
我们断言：任意非空开子叠 $U\subset\Bun_G$ 都含有一个基本（即半稳定）点。

先承认该断言来完成证明。需要看出 $\kappa$ 的纤维是连通分支。设 $b\in B(G)$ 是基本的，且 $U\subset\kappa^{-1}(\kappa(b))\subset\Bun_G$。则 $\Bun^b_G\subset U$，故 $\kappa^{-1}(\kappa(b))$ 必须连通，而 $b$ 是其中唯一的基本点。

现在证明该断言。取极小的 $b\in B(G)$ 使得 $\Bun^b_G\subset U$。由极小性，$\Bun^b_G\subset U$ 是开的。（否则 $U$ 内会出现泛化。）但 $U$ 作为 $\Bun_G$ 的开子叠上同调光滑、维数为 $0$，而 $\Bun^b_G$ 上同调光滑、维数为 $-\langle2\rho,\nu_b\rangle$。这意味着 $\langle2\rho,\nu_b\rangle=0$，故 $\nu_b$ 是中心的，即 $b$ 是基本的。
\end{proof}

\srcpage{85}
\begin{center}
\figP{085}
\end{center}

\subsection{与 $p$ 进群的光滑表示论的关系}

我们将进而讨论 $D_{\et}(\Bun_G,\Lambda)$ 与 $G_b(E)$ 的表示论之间的关系。

\begin{proposition}
对每个 $b\in B(G)$，有
\[
D_{\et}(\Bun^1_G,\Lambda)\cong\underbrace{D(G_b(E),\Lambda)}_{\substack{\text{$G_b(E)$ 在 $\Lambda$-模上的光滑表示}\\[1pt]\text{构成的阿贝范畴的导出范畴}}}.
\]
\end{proposition}

\subsubsection{第 0 步}

第 0 步是证明 $D_{\et}(*,\Lambda)\cong D(\Lambda)$。

这并非完全显然。为什么？容易证明：对几何点 $\Spa C$，有 $D_{\et}(\Spa C,\Lambda)\cong D(\Lambda)$，因为 $\Spa C$ 是上同调维数有限（事实上为 $0$）的空间菱形，故
\[
D_{\et}(\Spa C,\Lambda)\cong D((\Spa C)_{\et},\Lambda)=D(\Lambda),
\]
因为 $(\Spa C)_{\et}$ 是有限集的站点。

但 $*=\Spa\ol{\Fq}$ 不是菱形。要理解这里得到什么，必须分析沿 $\Spa C\to\Spa\ol{\Fq}$ 的下降。

\begin{proposition}
对任意小 v-叠 $X/\ol{\Fq}$，以及任意完备的代数闭非阿基米德域 $C$，拉回
\[
D_{\et}(X,\Lambda)\to D_{\et}(X\times_{\Spa\ol{\Fq}}\,\Spa C,\Lambda)
\]
是满忠实的。
\end{proposition}

由此可得
\[
D_{\et}(*,\Lambda)\hookrightarrow D_{\et}(\Spa C,\Lambda)\cong D(\Lambda).
\]
另一方面，你当然可以在任意空间上构造常值层，这给出交换图表
\[
\begin{tikzcd}
D_{\et}(*,\Lambda) \arrow[rr, shift left=0.7ex] & & D_{\et}(\Spa C,\Lambda) \arrow[ll, shift left=0.7ex] \\
& D(\Lambda) \arrow[ul] \arrow[ur] &
\end{tikzcd}
\]

\srcpage{86}

\subsubsection{第 1 步}

我们要证明
\[
D_{\et}([*/G_b(E)],\Lambda)\cong D(G_b(E),\Lambda),
\]
而且事实上对任意局部 pro-$p$ 群 $H$ 代替 $G_b(E)$ 也成立。类似地，我们要证明
\[
D_{\et}([\Spa C/G_b(E)],\Lambda)\cong D(G_b(E),\Lambda).
\]
想法是沿 $\Spa C\to[\Spa C/G_b(E)]$ 使用下降。需要小心，才能看出最终得到的是光滑表示的概念。

组织论证的更好方式是证明此时有
\[
D_{\et}([\Spa C/G_b(E)],\Lambda)\cong D([\Spa C/G_b(E)]_{\et},\Lambda)
\]
（尽管 $[\Spa C/G_b(E)]$ 不是菱形，而对菱形一般有这种比较）。现在，$[\Spa C/G_b(E)]_{\et}$ 等价于带有连续 $G_b(E)$-作用的集合构成的站点。（该陈述是说：连续作用映射 $G_b(E)\times S\to S$ 等价于通常意义下 $S$ 上的连续 $G_b(E)$-作用。）于是可以看出，该站点上的 $\Lambda$-模正是 $\Lambda$-模上的光滑 $G_b(E)$-表示（光滑指在离散 $\Lambda$-模上的连续作用）。

\subsubsection{第 2 步}

我们证明 $D_{\et}(\Bun^b_G,\Lambda)\cong D_{\et}([*/G_b(E)],\Lambda)$。研究映射
\[
[*/G_b(E)]\to\Bun^b_G.
\]
该映射上同调光滑，且纤维具有平凡上同调。

一般地，沿上同调光滑且纤维无环的纤维化作拉回是满忠实的，于是得到
\[
D_{\et}(\Bun^b_G,\Lambda)\hookrightarrow D_{\et}([*/G_b(E)],\Lambda).
\]
另一方面，该映射实际上是如下图表的分裂
\[
\begin{tikzcd}
& {[* / G_b(E)]} & \\
{[* / G_b(E)]} \arrow[rr] \arrow[ur, Rightarrow] & & \Bun^b_G \arrow[ul]
\end{tikzcd}
\]

\begin{corollary}[证明中的]
对任意完备的代数闭域 $C/\ol{\Fq}$，拉回诱导
\[
D_{\et}(\Bun_G,\Lambda)\cong D_{\et}(\Bun_G\times_{\ol{\Fq}}\,\Spa C,\Lambda).
\]
\end{corollary}

\begin{corollary}
$D_{\et}(\Bun_G,\Lambda)\cong D_{\et}(\Bun_G\times_{\ol{\Fq}}\,\Spa C,\Lambda)$，并且它容许一个无限半正交分解，其分量为
\[
D_{\et}(\Bun^b_G,\Lambda)\cong D(G_b(E),\Lambda).
\]
\end{corollary}

\begin{proof}
$\Bun_G$ 有分层，其分量为 $i^b\colon\Bun^b_G\hookrightarrow\Bun_G$。函子 $i^b_!$、$i^{b*}$ 诱导半正交分解。这是因为一般地，对开--闭分解
\[
i\colon Z\hookrightarrow X\leftarrow U\colon j
\]
总有半正交分解
\[
\begin{tikzcd}
D_{\et}(Z) \arrow[r, hook, "i_*"'] & D_{\et}(X) \arrow[l, bend left=40] \arrow[r, bend left=40] & D_{\et}(U) \arrow[l, hook', "j_!"']
\end{tikzcd}
\]

\srcpage{87}

不变性由形式论证推出。我们知道满忠实性 $D_{\et}(\Bun_G,\Lambda)\hookrightarrow D_{\et}(\Bun_G\times_{\ol{\Fq}}\,\Spa C,\Lambda)$，但也知道本质像包含 $i^b_!D_{\et}(\Bun^b_G\times_{\ol{\Fq}}\Spa C,\Lambda)$，于是由 dévissage 得到全部。
\end{proof}

\begin{remark}
分层之间如何相互作用，由上次课中的空间 $\pi_b\colon\mathcal{M}_b\to\Bun_G$ 编码。不过今天我们不展开这一点。
\end{remark}

\subsection{范畴 $D_{\et}(\Bun_G)$ 的结构}

回忆：若 $\mathcal{C}$ 是三角范畴，则 $X\in\mathcal{C}$ 称为紧的，是指 $\Hom_{\mathcal{C}}(X,-)$ 与一切直和交换。我们用 $\mathcal{C}^\omega$ 记紧对象的满子范畴。

\begin{theorem}
$D_{\et}(\Bun_G,\Lambda)$ 是紧生成的；复形 $A\in D_{\et}(\Bun_G,\Lambda)$ 紧，当且仅当对一切 $b\in B(G)$，$(i^b)^*A\in D_{\et}(\Bun^b_G,\Lambda)$ 紧，且除有限多个 $b$ 外都为 $0$。

等价地，它们位于由 $\mathrm{c}\text{-}\mathrm{Ind}^{G_b(E)}_K(\Lambda)$（$K\subset G_b(E)$ 为开 pro-$p$ 子群）生成的厚三角子范畴中。
\end{theorem}

\begin{warning}
$D_{\et}(\Bun_G,\Lambda)$ 中的紧对象不是 Verdier 自对偶的。表示论中已不陌生：$\mathrm{c}\text{-}\mathrm{Ind}^{G_b(E)}_K(\Lambda)$ 的对偶具有不可数维数。问题在于 $\mathrm{c}\text{-}\mathrm{Ind}^{G_b(E)}_K$ 不是容许的。
\end{warning}

\begin{theorem}
在紧对象子范畴 $D_{\et}(\Bun_G,\Lambda)^\omega\subset D_{\et}(\Bun_G,\Lambda)$ 上，存在 Bernstein--Zelevinsky 对偶函子
\[
\mathbb{D}_{\mathrm{BZ}}\colon(D_{\et}(\Bun_G,\Lambda)^\omega)^{\mathrm{op}}\to D_{\et}(\Bun_G,\Lambda)^\omega,
\]
使得 $\mathrm{RHom}(A,B)\cong\pi_\natural(\mathbb{D}_{\mathrm{BZ}}(A)\otimes^L_\Lambda B)$，其中 $\pi\colon\Bun_G\to*$，而 $\pi_\natural$ 是 $\pi^*$ 的左伴随。（它是 $R\pi_!$ 的一个扭转。）我们有 $\mathbb{D}_{\mathrm{BZ}}^2\cong\id$。

对基本的 $b\in B(G)$，$\mathbb{D}_{\mathrm{BZ}}$ 限制为 $D_{\et}(\Bun^b_G,\Lambda)\cong D(G_b(E),\Lambda)^\omega$ 上的自对偶，且与通常的 Bernstein--Zelevinsky 对偶一致。在对象层面，它发送
\[
\mathbb{D}_{\mathrm{BZ}}\bigl(\mathrm{c}\text{-}\mathrm{Ind}^{G_b(E)}_K\Lambda\bigr)=\mathrm{c}\text{-}\mathrm{Ind}^{G_b(E)}_K\Lambda.
\]
\end{theorem}

\subsection{ULA 层}

光滑对偶在容许表示上表现良好。因此我们将定义「$D_{\et}(\Bun_G)$ 版本的容许性」。对 $A\in D_{\et}(\Bun_G,\Lambda)$，有一个概念：$A$ 关于映射 $\Bun_G\to*$ 是普遍局部无环的（universally locally acyclic, ULA）\footnote{在通常的代数几何中，人们会期望每个（ind-）可构造层关于到 $*$ 的映射都是 ULA 的。}。这可以在分层上检验，且等价于：对一切 $b\in B(G)$，$(i^b)^*A\in D(G_b(E),\Lambda)$ 是容许的，即对一切开 pro-$p$ 子群 $K\subset G_b(E)$，
\[
[(i^b)^*A]^K\in D(\Lambda)
\]
是完美的（可由有限个有限投射 $\Lambda$-模构成的有限复形表示）。

这一类层在 Verdier 对偶下稳定，
\[
\mathbb{D}_{\Bun_G}(A)=\mathrm{RHom}(A,R\pi_!\Lambda),
\]
且满足 Verdier 双对偶：对任意 ULA 的 $A$，有 $A\xrightarrow{\sim}\mathbb{D}_{\Bun_G}(\mathbb{D}_{\Bun_G}(A))$。Verdier 对偶限制到分层上就是光滑对偶。

\begin{remark}
理想情况下，我们希望有 $\Bun_G$ 上「可构造复形」的概念；这些应当是紧对象，并且都关于 $\Bun_G\to*$ 普遍局部无环。然而这行不通！上述定理是最好的替代。

\srcpage{88}

但要注意紧不蕴涵 ULA，反之亦然。二者在表示论层面都已可见：$\mathrm{c}\text{-}\mathrm{Ind}^{G_b(E)}_K\Lambda$ 紧，但不是容许的（ULA）。反之，存在不是紧的 ULA 层。例如，考虑具有增长导子的超尖点表示的无限直和 $\bigoplus_{i=1}^\infty\pi_i$。它是容许的但不紧。事实上，对任意固定的开紧子群 $K$，有
\[
\Bigl(\bigoplus_{i=1}^\infty\pi_i\Bigr)^K=\bigoplus_{i=1}^{N(K)}\pi_i^K
\]
其中某个 $N(K)<\infty$，因为只有有限多个 $\pi_i$ 具有有界导子。

这种类型的表示正是你观察局部对称空间的上同调时看到的——它是容许的但不紧。
\end{remark}

\begin{warning}
在（局部）空间菱形上，有可构造复形的概念，它由小 v-叠上的下降给出。（它由 $j_!\Lambda$ 生成，其中 $j\colon U\to X$ 是 qcqs étale 映射。）这就是我们所说的可构造。但这是另一个不同的概念，而且本质上没有非零的 $A\in D_{\et}(\Bun_G,\Lambda)$ 在这个意义下可构造。
\end{warning}

\begin{example}
设 $X=\mathbb{D}_C$ 是闭单位圆盘，$i\colon\Spa C\hookrightarrow X$ 是原点的包含。那么你可能会期望 $i_*\Lambda$ 可构造，但它不是。问题在于补开集 $j\colon U=\mathbb{D}^*_C\hookrightarrow\mathbb{D}_C$ 不是拟紧的，因为 $\mathbb{D}^*_C$ 不是拟紧的。
\end{example}

事实上，可以证明：刚性解析簇 $X$ 上的可构造层在任何经典点的某个开邻域内局部常值。

要点：「有限生成」表示与「容许」表示的概念，连同 Bernstein--Zelevinsky 对偶与光滑对偶，都能推广到 $D_{\et}(\Bun_G,\Lambda)$。

\subsection{系数}

关于系数的注记：到目前为止，我们只允许 $\Lambda$ 满足 $n\Lambda=0$（某个与 $p$ 互素的 $n$）。理想情况下，我们希望 $\Lambda=\Ql$。出于与「可构造层」概念之所以棘手相同的原因，从 $\ZZ/\ell^n\ZZ$-系数过渡到 $\ZZ_\ell$-系数比通常更棘手。

我们可以定义
\[
D_{\et}(\Bun_G,\ZZ_\ell):=\varprojlim_n D_{\et}(\Bun_G,\ZZ/\ell^n\ZZ).
\]
但这与 $\ell$-进完备的 $\ZZ_\ell$-模上的表示相关\yizhu{紧接着的公式左端按字面为 $\ZZ/\ell^n\ZZ$-系数，与右端不相容；据上下文及下一式，左端疑应为 $D_{\et}(*,\ZZ_\ell)$。译文照原文排出。}{A（原文 p.88）}，例如
\[
D_{\et}(*,\ZZ/\ell^n\ZZ)=\varprojlim_n D(\ZZ/\ell^n\ZZ).
\]

但我们不想要 $\ell$-进完备向量空间上的表示；我们想要离散 $\ZZ_\ell$-向量空间上的表示。绕开这一点的通常技巧是考虑
\[
\Ind\Bigl(\varprojlim_n D_{\et}(\Bun_G,\ZZ/\ell^n\ZZ)^\omega\Bigr).
\]
在一点上这可行：
\[
\Ind\Bigl(\varprojlim_n D(\ZZ/\ell^n\ZZ)^\omega\Bigr)=D(\ZZ_\ell)
\]
是离散 $\ZZ_\ell$-模的导出 $\infty$-范畴，要点在于有限自由 $\ZZ_\ell$-模是 $\ell$-进完备的。

\srcpage{89}

然而这里行不通，因为紧对象不是容许的（不够「有限」）。

借助与 Dustin Clausen 合作发展的固体模思想，我们对任意 $\ZZ_\ell$-代数 $\Lambda$ 定义了 $D_{\et}(\Bun_G,\Lambda)$ 的一个版本，使得本讲的所有断言都成立。
